% !TEX program = xelatex
% 공통수학2 · Ⅰ. 도형의 방정식 · 2. 원의 방정식 · 01 원의 방정식 · 1/2차시
\documentclass[11pt,a4paper]{article}
\usepackage{kotex}
\usepackage{iftex}
\ifPDFTeX\else
  \setmainhangulfont{Noto Serif CJK KR}
  \setsanshangulfont{Noto Sans CJK KR}
\fi
\usepackage[margin=2.1cm]{geometry}
\usepackage{parskip}
\usepackage{enumitem}
\usepackage{array}
\usepackage{tabularx}
\usepackage{booktabs}
\usepackage{xcolor}
\usepackage{amssymb}
\usepackage{tikz}
\usepackage[most]{tcolorbox}
\usepackage{fancyhdr}
\usepackage{titlesec}

\definecolor{goldc}{HTML}{B07C22}
\definecolor{goldstrong}{HTML}{8A5F16}
\definecolor{indigoc}{HTML}{2B3B57}
\definecolor{indigosoft}{HTML}{6E82A6}
\definecolor{linec}{HTML}{D6DACB}
\definecolor{surface2}{HTML}{E4E8DC}
\definecolor{notebg}{HTML}{FBF3E3}
\definecolor{inksoft}{HTML}{52604F}

\titleformat{\section}{\normalfont\sffamily\bfseries\large\color{indigoc}}{\thesection}{0.6em}{}
\titlespacing{\section}{0pt}{16pt}{8pt}
\newtcolorbox{ideabox}{enhanced, breakable, colback=white, colframe=goldc, boxrule=0.5pt, leftrule=3pt, arc=2pt, left=10pt, right=10pt, top=8pt, bottom=8pt}
\newtcolorbox{calloutbox}{enhanced, breakable, colback=white, colframe=indigosoft, boxrule=0.5pt, leftrule=3pt, arc=2pt, left=10pt, right=10pt, top=7pt, bottom=7pt}
\newtcolorbox{notebox}{enhanced, breakable, colback=notebg, colframe=goldc, boxrule=0.5pt, leftrule=3pt, arc=2pt, left=10pt, right=10pt, top=7pt, bottom=7pt}
\newtcolorbox{answerbox}{enhanced, breakable, colback=white, colframe=linec, boxrule=0.5pt, arc=2pt, left=10pt, right=10pt, top=7pt, bottom=7pt}
\newcommand{\lessonbanner}[3]{%
  \begin{tcolorbox}[enhanced, colback=indigoc, colframe=indigoc, boxrule=0pt, arc=0pt,
    left=14pt, right=14pt, top=14pt, bottom=10pt, borderline south={2.4pt}{0pt}{goldc}, fontupper=\color{white}]
    {\sffamily\footnotesize\color{white!85!black} #1}\\[4pt]{\sffamily\bfseries\Large #2}\\[3pt]{\sffamily\small\color{white!88!black} #3}
  \end{tcolorbox}}
\pagestyle{fancy}\fancyhf{}\renewcommand{\headrulewidth}{0pt}
\fancyfoot[L]{\footnotesize\color{inksoft} 공통수학2 · 2026학년도 2학기 · 지평선고등학교}
\fancyfoot[R]{\footnotesize\color{inksoft} 교사용 지도안 -- 배포 금지}

\begin{document}
\lessonbanner{공통수학2 · Ⅰ. 도형의 방정식 · 2. 원의 방정식}{01 원의 방정식 --- 1/2차시}{원의 방정식(표준형) · 교사용 지도안 (2026학년도 2학기)}
\vspace{10pt}

\noindent\begin{tabularx}{\linewidth}{@{}>{\bfseries\color{indigoc}}p{2.6cm}X@{}}
\toprule
대상 & 고1 · 2개 반 · 반당 13명 \\
차시 & 50분 (도입 10 · 전개 30 · 정리 10) \\
성취기준 & 원의 방정식을 구하고, 그래프를 그릴 수 있다. \\
선행 학습 & 1중단원 --- 두 점 사이의 거리, 내분점 \\
\bottomrule
\end{tabularx}

\section{학습 목표 \& Big Idea}
\begin{ideabox}
\textbf{\color{goldstrong}영속적 이해 (Big Idea)}\\
원은 `중심으로부터 같은 거리에 있는 점들의 모임'이라는 정의 그 자체가 방정식이 된다. 도형의 정의를 그대로 식으로 옮기면 방정식을 얻는다.
\vspace{4pt}
\small\color{inksoft} 핵심개념: \textbf{정의의 번역} \quad\textbullet\quad 핵심개념: \textbf{일정한 거리} \quad\textbullet\quad 일반화: \textbf{도형의 정의는 곧 그 도형의 방정식이다}
\end{ideabox}
\begin{calloutbox}
\textbf{차시 학습 목표}
\begin{itemize}[leftmargin=1.4em, itemsep=2pt, topsep=4pt]
  \item 원의 정의로부터 원의 방정식 $(x-a)^2+(y-b)^2=r^2$을 유도할 수 있다.
  \item 중심과 반지름, 또는 지름의 양 끝 점이 주어졌을 때 원의 방정식을 구할 수 있다.
\end{itemize}
\end{calloutbox}

\section{질문의 연결}
\begin{tabularx}{\linewidth}{@{}>{\bfseries\color{indigoc}\small}p{2.4cm}X@{}}
사실적 질문 & 중심 $(a,b)$, 반지름 $r$인 원 위의 점 $(x,y)$는 어떤 식을 만족할까? \\[6pt]
개념적 질문 & 원의 `정의'와 원의 `방정식'은 어떻게 같은 것을 다르게 표현한 것일까? \\[6pt]
논쟁적 질문 & 원형 관개 농업처럼, 수학의 원이 실제 세계의 효율적인 설계에 쓰이는 다른 예는 무엇이 있을까? \\
\end{tabularx}

\section{수업 흐름}
\begin{tabularx}{\linewidth}{@{}>{\centering\arraybackslash}p{1.6cm}X>{\raggedright\arraybackslash\small}p{3.1cm}@{}}
\toprule
\textbf{단계} & \textbf{활동 \& 발문} & \textbf{ATL / VTR} \\
\midrule
\textcolor{goldstrong}{\textbf{도입}}\newline\footnotesize 10분 &
\textbf{마음공부 (2분)} --- 유무념 대조.\newline
\textbf{생각 열기 (8분)} --- 원형 관개 농업 사진 제시(반지름을 따라 회전하는 스프링클러) $\to$ 반지름 $10$인 정부 상징 문양을 원점 중심으로 좌표평면에 나타내고, $\overline{\text{OP}}=10$으로부터 $x^2+y^2=$ 얼마인지 채워보게 한다. &
ATL: 사고(정의 $\to$ 식 번역)\newline VTR: See-Think-Wonder \\
\addlinespace
\textcolor{goldstrong}{\textbf{전개}}\newline\footnotesize 30분 &
\textbf{개념 형성 (12분)} --- 중심 $C(a,b)$, 반지름 $r$인 원 위의 점 $P(x,y)$에 대해 $\overline{\text{CP}}=r$(두 점 사이 거리 공식)로부터 $(x-a)^2+(y-b)^2=r^2$ 유도 $\to$ 문제 1.\newline
\textbf{적용 (18분)} --- 지름의 양 끝점이 주어질 때 중심은 중점, 반지름은 중점까지 거리임을 확인(예제 1) $\to$ 문제 2 짝 활동 $\to$ 좌표축에 접하는 원(지도자료) 간단 소개. &
ATT: 촉진자 역할\newline ATL: 사고·적용 \\
\addlinespace
\textcolor{goldstrong}{\textbf{정리}}\newline\footnotesize 10분 &
\textbf{개념 연결 질문 (5분)} --- ``원의 정의에서 `일정한 거리'가 방정식의 어느 부분이 되었는가?''\newline
\textbf{성찰 (5분)} --- 감사 일기 1문장 + 다음 차시 예고(일반형, 세 점을 지나는 원). &
마음공부 · 감사 일기 \\
\bottomrule
\end{tabularx}

\begin{calloutbox}
\textbf{지도상의 유의점} --- 공식을 암기시키기 전에 반드시 `두 점 사이의 거리 공식 $\to$ 양변 제곱'의 유도 과정을 학생이 직접 손으로 쓰게 한다. 지름의 양 끝점 문제는 `중점 공식(반지름) + 중점 공식(중심)'을 순서대로 적용하는 2단계 사고임을 강조한다.
\end{calloutbox}

\section{형성평가 \& 답안}
\begin{minipage}[t]{0.48\linewidth}
\begin{answerbox}
\textbf{문제 1}\\
\small ⑴ 중심$(-3,0)$, 반지름 $2$ \quad ⑵ 중심 원점, 점$(3,-4)$ 지남\\[4pt]
\textbf{\color{goldstrong}답} \; ⑴ $(x+3)^2+y^2=4$ \quad ⑵ $x^2+y^2=25$
\end{answerbox}
\end{minipage}\hfill
\begin{minipage}[t]{0.48\linewidth}
\begin{answerbox}
\textbf{문제 2 (예제 1과 동일 유형)}\\
\small A$(-4,10)$, B$(8,-6)$을 지름의 양 끝점\\[4pt]
\textbf{\color{goldstrong}답} \; 중심$(2,2)$, 반지름 $10$ $\to$ $(x-2)^2+(y-2)^2=100$
\end{answerbox}
\end{minipage}

\section{성취수준 (이 차시 범위)}
\begin{tabularx}{\linewidth}{@{}>{\bfseries\color{goldstrong}}p{0.9cm}X@{}}
\toprule
A & 원의 정의로부터 원의 방정식을 유도할 수 있으며, 주어진 조건으로 원의 방정식을 구하고 그래프를 그릴 수 있다. \\
B & 원의 방정식을 구하고 그래프를 그릴 수 있다. \\
C & 중심과 반지름이 주어졌을 때 원의 방정식을 구하고 그래프를 그릴 수 있다. \\
D & 표준형 원의 방정식의 그래프를 그릴 수 있다. \\
E & 안내된 절차에 따라 원의 방정식의 그래프를 그릴 수 있다. \\
\bottomrule
\end{tabularx}

\section{다음 차시}
\begin{calloutbox}
\textbf{01 원의 방정식 --- 2/2차시.} 세 점을 지나는 원의 방정식(외접원)과, 이차방정식 $x^2+y^2+Ax+By+C=0$이 나타내는 도형(일반형)을 다룬다.
\end{calloutbox}
\end{document}
